%! Author = denis
%! Date = 11/08/2026

\begin{frame}{Identidades lógicas: Lei da Anulação}
    \begin{block}{Definição}
        \begin{itemize}
            \item uma operação \textbf{OR} com $1$ sempre resulta em $1$;
            \item uma operação \textbf{AND} com $0$ sempre resulta em $0$.
        \end{itemize}
    \end{block}

    \vspace{0.3cm}

    \begin{columns}[T]

        % Coluna OR
        \begin{column}{0.5\textwidth}
            \begin{block}{Anulação OR}
                \[
                    \boxed{A + 1 = 1}
                \]
                \[
                    0 + 1 = 1
                \]
                \[
                    1 + 1 = 1
                \]

                Portanto, a saída será sempre $1$.
            \end{block}
        \end{column}

        % Coluna AND
        \begin{column}{0.5\textwidth}
            \begin{block}{Anulação AND}
                \[
                    \boxed{A \cdot 0 = 0}
                \]
                \[
                    0 \cdot 0 = 0
                \]
                \[
                    1 \cdot 0 = 0
                \]

                Portanto, a saída será sempre $0$.
            \end{block}
        \end{column}

    \end{columns}


    \begin{alertblock}{Interpretação}
        O valor $1$ é dominante na operação \textbf{OR}, enquanto
        o valor $0$ é dominante na operação \textbf{AND},
        independentemente do valor de $A$.
    \end{alertblock}
    
\end{frame}